Transformationsvektoren (Um einen Vektor in anderen Koordinaten darzustellen):\\
\begin{tabular}{c | c  l}
             & $(x, y, [z])^\top=F(u)=$                                                                                    &                                                                            \\ \midrule
    Polar    & $\begin{pmatrix} \cos(\varphi) & -r\sin(\varphi) \\ \sin(\varphi) & r\cos(\varphi) \end{pmatrix}$           & $0 \le \varphi < 2\pi$                                                     \\
    Zylinder & $\begin{pmatrix} r\cos (\varphi) \\ r\sin (\varphi) \\ z \end{pmatrix}$                                     & $0 \le \varphi < 2 \pi$                                                    \\
    Kugel    & $\begin{pmatrix} r\cos(\varphi) \sin(\theta) \\ r\sin(\varphi) \sin(\theta) \\ r\cos(\theta) \end{pmatrix}$ & $\begin{matrix}0 \le \varphi < 2 \pi \\ 0 \le \theta \le \pi \end{matrix}$
\end{tabular}

\subsection{Transformationen}
Koordinaten $x=(x,y,z,...) \in \mathbb R^n$ in karthesischer Darstellung werden durch eine Abbildung $F:\mathbb R^n \rightarrow \mathbb R^n, x=F(u)$ dargestellt, wobei $u$ die Koordinaten in anderer Basis sind (z.B. Kugel)
\begin{description}
    \item[regulär] Punkt, in welchem $\det_F(u)\ne 0$
    \item[lokale Basisvektoren] Spalten $c_1,...,c_n$ von $J_F(u)$, abhängig vom Punkt $u$ bzw. $x$
    \item[Maßfaktor] Länge der Basisvektoren: $h_i=||c_i||$
    \item[orthogonale Koordinaten] Basisvektoren sind alle orthogonal zueinander $\Rightarrow B=\begin{pmatrix}\frac{1}{h_1}c_1 & \dots & \frac{1}{h_n}c_n\end{pmatrix}$ ist Orthogonalmatrix ($B^\top=B^{-1}$)
\end{description}

\subsection{Vektorfelder in Krummlinigen Koordinaten}
$G:\mathbb R^n \rightarrow \mathbb R^n$ Vektorfeld bzgl. karth. Koordinaten \\
\textbf{I. Koordinatentransformation im Definitionsraum} \\
\begin{equation*}
    u=F(x) \rightarrow \tilde{G}(u)=\begin{pmatrix}
        \tilde{g_1}(u) \\ \vdots \\ \tilde{g_n}(u)
    \end{pmatrix}=G(F(u))
\end{equation*}
Darstellung weiterhin in karthesischen Koordinaten: \\
$\tilde{G}(u)=\tilde{g_1}(u)e_1 + \dots + \tilde{g_n}(u)e_n$

\textbf{II. Koordinatentransformation im Bildraum} \\
\begin{equation*}
    \tilde{G}(u)=\underbrace{\begin{pmatrix}
            |       &     & |       \\
            e_{u_1} & ... & e_{u_n} \\
            |       &     & |
        \end{pmatrix}}_{B}\hat{G}(u) \rightarrow \hat{G}(u)=B^{-1}\tilde{G}(u)
\end{equation*}
Darstellung jetzt in neuen Koordinaten: \\
$\hat{G}(u)=\hat{g_1}(u)e_{u_1} + \dots + \hat{g_n}(u)e_{u_n}$